\chapter{Rectifiable parts and projection tests}\label{sec:rectifiable-projection-structure}
A set can contain both smooth pieces and a part that meets every smooth piece in zero measure. This section develops the measurable decomposition and elementary projection tests in \cite{Sim18}, Chapter 3, Definition 3.1 and Lemmas 3.2--3.3, printed pp.~94--95. We assume $1\leq m<n$. The almost-everywhere projection theorem of Besicovitch and Federer is a deeper result and is not used in these proofs.

\begin{definition}[Pure unrectifiability]\label{def:pure-unrectifiable-set}
An $\mathcal H^m$-measurable set $P\subset\mathbb R^n$ is purely $m$-unrectifiable if $\mathcal H^m(P\cap R)=0$ for every measurable countably $m$-rectifiable set $R$. Equivalently, $P$ contains no measurable countably $m$-rectifiable subset of positive $\mathcal H^m$ measure. The equivalence uses the fact that measurable subsets of a countably rectifiable set are countably rectifiable, directly from Definition \ref{def:rectifiable-set}. This is the measurable-set convention for \cite{Sim18}, Definition 3.1.
\end{definition}

\begin{lemma}[Borel cores for a sigma-finite Hausdorff measure]\label{lem:hausdorff-borel-core}
If $A\subset\mathbb R^n$ is $\mathcal H^m$-measurable and sigma-finite for that measure, then there are Borel sets $C\subset A\subset B$ with $\mathcal H^m(B\setminus C)=0$. If $A$ is countably $m$-rectifiable, both $B$ and $C$ are countably $m$-rectifiable as well.
\end{lemma}
\begin{proof}
Write $A$ as a countable union of measurable finite-measure sets $A_j$. Borel regularity gives a Borel hull $B_j\supset A_j$ with the same finite measure. Since $A_j$ is measurable, $\mathcal H^m(B_j\setminus A_j)=0$. Apply Borel regularity again to obtain a Borel null set $Z_j$ containing $B_j\setminus A_j$. Then $C_j=B_j\setminus Z_j$ is Borel, lies in $A_j$, and differs from it by a null set. Put $C=\bigcup_j C_j$ and $B=\bigcup_j B_j$. These have the required inclusions and null difference. A subset of a countably rectifiable set is covered by the same Lipschitz images up to the inherited null set; adding a Hausdorff-null set also preserves rectifiability. Thus the last assertion follows from the definition. This argument does not require local finiteness of the unweighted area of $A$.
\end{proof}

\begin{proposition}[Testing pure unrectifiability on smooth patches]\label{prop:pure-unrectifiable-patch-test}
An $\mathcal H^m$-measurable set $P$ is purely $m$-unrectifiable if and only if $\mathcal H^m(P\cap N)=0$ for every finite-area embedded $C^1$ graph patch $N$ of dimension $m$.
\end{proposition}
\begin{proof}
Every such patch is countably rectifiable, so the forward implication is part of Definition \ref{def:pure-unrectifiable-set}. Conversely let $R$ be countably rectifiable. Theorem \ref{thm:rectifiable-c1-decomposition} and Lemma \ref{lem:finite-area-c1-patches} cover $R$, apart from a Hausdorff-null subset, by countably many finite-area smooth patches $N_j$. Their intersections with $P$ are null by assumption. Their countable union and the exceptional null set still have zero measure, proving $\mathcal H^m(P\cap R)=0$. Local finiteness of $P$ is not needed for this implication.
\end{proof}

\begin{proposition}[A maximal rectifiable part of a finite-measure set]\label{prop:finite-maximal-rectifiable-part}
Let $A$ be a Borel set with $\mathcal H^m(A)<\infty$. There exists a Borel countably $m$-rectifiable subset $R\subset A$ such that $P=A\setminus R$ is purely $m$-unrectifiable. Its measure is
\[
 \mathcal H^m(R)=\sup\{\mathcal H^m(S):S\subset A\text{ Borel and countably }m\text{-rectifiable}\}.
\]
This is the finite-measure step in \cite{Sim18}, Lemma 3.2.
\end{proposition}
\begin{proof}
Call the supremum $\alpha\in[0,\mathcal H^m(A)]$. Choose Borel rectifiable $S_j\subset A$ with $\mathcal H^m(S_j)>\alpha-1/j$, allowing $S_j=\varnothing$ when $\alpha=0$. Their union $R$ is Borel and countably rectifiable. It contains every $S_j$, so $\mathcal H^m(R)\geq\alpha$; by its own admissibility it has measure at most $\alpha$. Hence equality holds.

If the residual set contained a measurable rectifiable subset $E$ of positive measure, Lemma \ref{lem:hausdorff-borel-core} would give a Borel rectifiable core $E'\subset E$ of the same positive measure. The disjoint union $R\cup E'$ would be an admissible Borel rectifiable subset of $A$ with measure greater than $\alpha$, a contradiction. By Definition \ref{def:pure-unrectifiable-set}, the residual is purely unrectifiable. Finiteness of $\alpha$ is the reason for doing this step on finite-measure pieces.
\end{proof}

\begin{theorem}[Decomposition into rectifiable and purely unrectifiable parts]\label{thm:rectifiable-pure-decomposition}
Every $\mathcal H^m$-measurable sigma-finite set $A\subset\mathbb R^n$ admits a disjoint decomposition $A=R\cup P$, where $R$ is Borel and countably $m$-rectifiable, and $P$ is measurable and purely $m$-unrectifiable. The decomposition is unique modulo $\mathcal H^m$-null sets. This is the measurable form of \cite{Sim18}, Lemma 3.2, with uniqueness made explicit.
\end{theorem}
\begin{proof}
Use Lemma \ref{lem:hausdorff-borel-core} to choose a Borel core of $A$ with null remainder. Partition that core into countably many Borel sets $A_j$ of finite measure. To see this partition exists, first use finite-measure Borel hulls of a sigma-finite measurable cover and intersect them with the core, then remove preceding sets. Apply Proposition \ref{prop:finite-maximal-rectifiable-part} to each $A_j$, obtaining $A_j=R_j\cup P_j$. Set $R=\bigcup_jR_j$ and $P=A\setminus R$. The union $R$ is Borel and rectifiable. For every rectifiable set $S$, each $S\cap P_j$ is null, and the omitted core remainder is null. Thus $S\cap P$ is null, proving pure unrectifiability.

For uniqueness, suppose also $A=\widetilde R\cup\widetilde P$ with the same properties. The measurable set $R\setminus\widetilde R$ lies in $R\cap\widetilde P$, which is null because $R$ is rectifiable and $\widetilde P$ is purely unrectifiable. Interchanging the decompositions gives $\mathcal H^m(\widetilde R\setminus R)=0$. Hence $R\mathbin\triangle\widetilde R$ and consequently $P\mathbin\triangle\widetilde P$ are null. Uniqueness is a measure assertion and does not assert equality of topological supports.
\end{proof}

\begin{corollary}[The rectifiable part captures every rectifiable subset]\label{cor:rectifiable-part-maximal-null}
For the decomposition in Theorem \ref{thm:rectifiable-pure-decomposition}, every measurable countably $m$-rectifiable subset $S\subset A$ satisfies $\mathcal H^m(S\setminus R)=0$. More generally, the rectifiable part of a measurable subset $B\subset A$ is $B\cap R$ modulo a null set.
\end{corollary}
\begin{proof}
The first difference is contained in $S\cap P$, which is null by pure unrectifiability. For the second assertion, $B\cap R$ is rectifiable, and $B\cap P$ remains purely unrectifiable by the definition. The two form a decomposition of $B$. Choosing a Borel core of $B\cap R$ changes it only by a null set, so the uniqueness statement of Theorem \ref{thm:rectifiable-pure-decomposition} applies. Thus restriction commutes with the decomposition at the measure level, even if the chosen set representatives change.
\end{proof}

\begin{lemma}[A regular coordinate projection detects positive area locally]\label{lem:coordinate-projection-detects-area}
Let $N$ be an embedded $m$-dimensional $C^1$ manifold in $\mathbb R^n$, and let $E\subset N$ be measurable with positive $\mathcal H^m$ measure. Then there are $m$ coordinate axes and a graph patch $N_0\subset N$ such that the corresponding orthogonal projection $p_I$ is injective on $N_0$, has a Lipschitz inverse on its image, and
\[
 \mathcal H^m(p_I(E\cap N_0))>0.
\]
This is the local inverse-function argument in \cite{Sim18}, Lemma 3.3, p.~95.
\end{lemma}
\begin{proof}
A regular chart has an $n\times m$ derivative matrix of rank $m$, so at every point some $m\times m$ coordinate minor is invertible. The inverse function theorem writes a neighborhood in $N$ as a graph over those coordinate axes. Shrink to an open base ball with compact closure in the graph domain. Its graph inverse has bounded derivative and is Lipschitz on that convex ball. These neighborhoods give a countable cover of $N$, allowing different coordinate projections for different neighborhoods. Since $E$ has positive measure, one covered piece $E\cap N_0$ has positive measure. If its projection were null, the Hausdorff-measure Lipschitz bound for the inverse graph map would make the piece itself null, a contradiction. Measurability of the projected set follows by taking a Borel representative of the piece and a null remainder: the graph projection is a homeomorphism on its patch and maps the remainder to a null set. No global injectivity of the projection on all of $N$ is claimed.
\end{proof}

\begin{theorem}[Null coordinate projections imply pure unrectifiability]\label{thm:null-coordinate-projections-pure}
Let $S\subset\mathbb R^n$ be $\mathcal H^m$-measurable. If
\[
 \mathcal H^m(p_I(S))=0
 \quad\text{for every }I\subset\{1,\ldots,n\},\quad |I|=m,
\]
then $S$ is purely $m$-unrectifiable. Here $p_I$ is projection onto the coordinate $m$-plane indexed by $I$, and the zero-measure assertion may be read in the outer-measure sense. This is \cite{Sim18}, Lemma 3.3.
\end{theorem}
\begin{proof}
If $S$ were not purely unrectifiable, Proposition \ref{prop:pure-unrectifiable-patch-test} would supply a finite-area smooth patch $N$ with $\mathcal H^m(S\cap N)>0$. Lemma \ref{lem:coordinate-projection-detects-area}, applied to that common part, would then give a coordinate projection with positive projected measure. Its image lies inside $p_I(S)$, contradicting the assumed null outer measure. There are only finitely many coordinate $m$-planes, but the assumption must hold for all of them. A single null projection does not suffice, as an $m$-plane with a rank-deficient projection can have that property.
\end{proof}

\begin{proposition}[Coordinate projection Jacobians control rectifiable area with multiplicity]\label{prop:coordinate-projection-area-bounds}
Let $M\subset\mathbb R^n$ be measurable, countably $m$-rectifiable and locally finite in $\mathcal H^m$ measure, and put $q=\binom nm$. For each coordinate $m$-plane let $N_I(y)=\mathcal H^0(M\cap p_I^{-1}(y))$. Then
\[
 \mathcal H^m(M)\leq\sum_{|I|=m}\int_{p_I(\mathbb R^n)}N_I(y)\,d\mathcal H^m(y)
 \leq\sqrt q\,\mathcal H^m(M).
\]
Infinite values are allowed. This is an area-formula refinement of the projection test; the integrals count multiplicity rather than only image measure.
\end{proposition}
\begin{proof}
At a point with approximate tangent choose an orthonormal $n\times m$ frame $Q$ for it. The area Jacobian $J_I$ of $p_I$ restricted to this tangent is the absolute value of the $I$-minor of $Q$. The Cauchy--Binet identity gives $\sum_IJ_I^2=\det(Q^*Q)=1$. For this finite nonnegative vector, its sum lies between its Euclidean norm and $\sqrt q$ times that norm, hence $1\leq\sum_IJ_I\leq\sqrt q$. Integrate on $M$ and use Theorem \ref{thm:rectifiable-weighted-area-local} with $f=p_I$ and weight one. That theorem identifies $\int_MJ_I$ with the multiplicity integral in the statement. The finite sum commutes with the integral, proving both bounds. If every multiplicity is at most an integer $N$ almost everywhere, the first bound in particular gives $\mathcal H^m(M)\leq N\sum_I\mathcal H^m(p_I(M))$.
\end{proof}
